\documentclass[10pt,aspectratio=169]{beamer}
\usepackage[utf8]{inputenc}
\usepackage[T1]{fontenc}
\usepackage[french]{babel}
\usepackage{amssymb}
\usepackage{booktabs}
\usepackage{graphicx}
\usepackage{inconsolata}
\usepackage{listings}
\usepackage{microtype}
\usepackage{tikz}
\usepackage{xcolor}
\usetikzlibrary{arrows.meta,positioning}

\definecolor{ink}{HTML}{202124}
\definecolor{muted}{HTML}{5F6368}
\definecolor{line}{HTML}{DADCE0}
\definecolor{paper}{HTML}{F8F9FA}
\definecolor{accent}{HTML}{8A1538}
\definecolor{green}{HTML}{0B6B57}
\definecolor{blue}{HTML}{245B8E}
\definecolor{amber}{HTML}{9A5B00}

\setbeamercolor{normal text}{fg=ink,bg=white}
\setbeamercolor{structure}{fg=accent}
\setbeamercolor{frametitle}{fg=ink,bg=white}
\setbeamercolor{title}{fg=ink}
\setbeamercolor{subtitle}{fg=muted}
\setbeamercolor{block title}{fg=white,bg=accent}
\setbeamercolor{block body}{fg=ink,bg=paper}
\setbeamercolor{alerted text}{fg=accent}

\setbeamertemplate{navigation symbols}{}
\setbeamertemplate{itemize item}{\color{accent}\small$\blacktriangleright$}
\setbeamertemplate{itemize subitem}{\color{green}\scriptsize$\bullet$}
\setbeamertemplate{frametitle}{
  \vspace{0.15cm}
  \insertframetitle\par
  \vspace{0.1cm}
  {\color{line}\rule{\linewidth}{0.45pt}}
}
\setbeamertemplate{footline}{
  \leavevmode%
  \hbox{%
    \begin{beamercolorbox}[wd=.76\paperwidth,ht=2.5ex,dp=1ex,leftskip=0.35cm]{author in head/foot}%
      Python pour économistes -- L2 Économie
    \end{beamercolorbox}%
    \begin{beamercolorbox}[wd=.24\paperwidth,ht=2.5ex,dp=1ex,rightskip=0.35cm plus1fil]{date in head/foot}%
      \hfill\insertframenumber/\inserttotalframenumber
    \end{beamercolorbox}%
  }%
}

\setbeamersize{text margin left=0.65cm,text margin right=0.65cm}


\lstdefinelanguage{terminal}{
  morekeywords={ls,cd,pwd,mkdir,python,python3,code,dir,type,echo},
  sensitive=true
}

\lstset{
  basicstyle=\ttfamily\small,
  backgroundcolor=\color{paper},
  frame=single,
  framerule=0.35pt,
  rulecolor=\color{line},
  keywordstyle=\color{accent}\bfseries,
  commentstyle=\color{muted},
  stringstyle=\color{green},
  showstringspaces=false,
  breaklines=true,
  columns=fullflexible,
  keepspaces=true,
  upquote=true,
  language=Python,
  literate=
   {à}{{\`a}}1 {â}{{\^a}}1 {ä}{{\"a}}1
   {ç}{{\c c}}1
   {é}{{\'e}}1 {è}{{\`e}}1 {ê}{{\^e}}1 {ë}{{\"e}}1
   {î}{{\^\i}}1 {ï}{{\"\i}}1
   {ô}{{\^o}}1 {ö}{{\"o}}1
   {ù}{{\`u}}1 {û}{{\^u}}1 {ü}{{\"u}}1
   {À}{{\`A}}1 {Â}{{\^A}}1 {Ä}{{\"A}}1
   {Ç}{{\c C}}1
   {É}{{\'E}}1 {È}{{\`E}}1 {Ê}{{\^E}}1 {Ë}{{\"E}}1
   {Î}{{\^I}}1 {Ï}{{\"I}}1
   {Ô}{{\^O}}1 {Ö}{{\"O}}1
   {Ù}{{\`U}}1 {Û}{{\^U}}1 {Ü}{{\"U}}1
   {œ}{{\oe}}1 {Œ}{{\OE}}1
   {–}{{--}}1 {—}{{---}}1
}

\newcommand{\code}[1]{\texttt{#1}}
\newcommand{\lesson}[1]{\textbf{\textcolor{accent}{#1}}}
\newcommand{\muted}[1]{\textcolor{muted}{#1}}


\title{Cours 2 -- Décrire, comparer, visualiser des données OWID}
\subtitle{Moyenne, médiane, dispersion et premiers pas avec pandas}
\author{Etienne Dagorn}
\institute{Université de Lille -- L2 Économie}
\date{2026--2027}

\begin{document}

\begin{frame}
  \titlepage
\end{frame}

\begin{frame}{Objectif de la séance}
\begin{block}{But}
Passer de petits objets Python à une première analyse descriptive : lire des données, calculer quelques indicateurs, visualiser et interpréter avec prudence.
\end{block}
\begin{itemize}
  \item Python reste l'objectif principal.
  \item Les statistiques servent de langage pour résumer les données.
  \item Les exemples suivent la logique des TPs OWID.
\end{itemize}
\end{frame}

\begin{frame}{Où va le cours ?}
\begin{tabular}{@{}lll@{}}
\toprule
Séance & Python & Statistique support \\
\midrule
TP1 & variables, listes & moyenne, min, max, pourcentage \\
TP2 & boucles, conditions & proportions, seuils \\
TP3 & fonctions & comparaison de groupes \\
TP4 & graphiques, simulation & distribution, stabilisation d'une moyenne \\
TP5 & \code{pandas} & mini-projet descriptif OWID \\
\bottomrule
\end{tabular}
\vspace{0.45cm}
\begin{block}{Fil rouge}
Une bonne analyse décrit une distribution, pas seulement un chiffre isolé.
\end{block}
\end{frame}

\begin{frame}{Donnée, observation, variable}
\begin{tabular}{@{}lll@{}}
\toprule
Mot & Sens simple & Exemple OWID \\
\midrule
Observation & une ligne du tableau & France en 2021 \\
Variable & une colonne du tableau & PIB par habitant \\
Valeur & une case du tableau & 46000 \\
Unité & sens de la valeur & dollars par habitant \\
\bottomrule
\end{tabular}
\vspace{0.45cm}
\begin{block}{Réflexe}
Avant de calculer : quelle est l'observation ? quelle est la variable ? quelle est l'unité ?
\end{block}
\end{frame}

\begin{frame}[fragile]{Même information, trois formes}
\begin{columns}[T,totalwidth=\textwidth]
\begin{column}{0.32\textwidth}
\begin{lstlisting}
valeurs = [46, 54, 34]
\end{lstlisting}
\end{column}
\begin{column}{0.32\textwidth}
\begin{lstlisting}
gdp = {
    "France": 46,
    "Germany": 54,
    "Spain": 34
}
\end{lstlisting}
\end{column}
\begin{column}{0.32\textwidth}
\begin{lstlisting}
donnees = [
    {"pays": "France", "gdp": 46},
    {"pays": "Germany", "gdp": 54},
    {"pays": "Spain", "gdp": 34}
]
\end{lstlisting}
\end{column}
\end{columns}
\begin{block}{Idée}
Plus les données deviennent riches, plus une table devient pratique. C'est le rôle d'un DataFrame \code{pandas}.
\end{block}
\end{frame}

\begin{frame}{Langage, bibliothèque, environnement}
\begin{tabular}{@{}lll@{}}
\toprule
Terme & Rôle & Exemple \\
\midrule
Langage & règles pour écrire des programmes & Python \\
Interpréteur & programme qui lit le langage & \code{python3} \\
Bibliothèque & outils ajoutés au langage & \code{pandas}, \code{matplotlib} \\
Environnement & ensemble installé sur la machine & Python + packages \\
\bottomrule
\end{tabular}
\vspace{0.45cm}
\begin{block}{Erreur fréquente}
Si \code{import pandas} échoue, Python fonctionne peut-être très bien : c'est la bibliothèque ou l'environnement qui manque.
\end{block}
\end{frame}

\begin{frame}[fragile]{Pseudo-code : décrire une liste}
\begin{lstlisting}
# Objectif : resumer une liste de valeurs economiques
# Entree : une liste de nombres
# 1. Compter les observations
# 2. Calculer la moyenne
# 3. Trier pour trouver la mediane
# 4. Reperer minimum et maximum
# 5. Observer la dispersion
# Sortie : chiffres + interpretation courte
\end{lstlisting}
\begin{block}{Programmation}
Le pseudo-code dit quoi faire. Python impose ensuite une syntaxe précise pour le faire.
\end{block}
\end{frame}

\begin{frame}{Moyenne et médiane : ce qu'elles disent}
\begin{tabular}{@{}lll@{}}
\toprule
Indicateur & Question & Limite \\
\midrule
Moyenne & niveau général & sensible aux valeurs extrêmes \\
Médiane & valeur centrale & ignore l'ampleur des écarts \\
Minimum & plus petite valeur & ne résume pas le groupe \\
Maximum & plus grande valeur & peut être atypique \\
\bottomrule
\end{tabular}
\vspace{0.45cm}
\begin{block}{À retenir}
Moyenne et médiane résument le centre. Elles ne suffisent pas à décrire la forme de la distribution.
\end{block}
\end{frame}

\begin{frame}[fragile]{Même moyenne, même médiane, distributions différentes}
\begin{lstlisting}
groupe_a = [8, 9, 10, 11, 12]
groupe_b = [0, 5, 10, 15, 20]
\end{lstlisting}
\begin{tabular}{@{}llll@{}}
\toprule
Groupe & Moyenne & Médiane & Lecture \\
\midrule
A & 10 & 10 & valeurs concentrées \\
B & 10 & 10 & valeurs très dispersées \\
\bottomrule
\end{tabular}
\vspace{0.4cm}
\begin{block}{Question}
Si ces valeurs sont des revenus, les deux groupes racontent-ils la même situation économique ?
\end{block}
\end{frame}

\begin{frame}{Distribution : regarder toutes les valeurs}
\begin{itemize}
  \item Une distribution décrit la façon dont les valeurs se répartissent.
  \item On peut regarder le centre : moyenne, médiane.
  \item On peut regarder les extrêmes : min, max, étendue.
  \item On peut regarder la dispersion : écarts à la moyenne, variance, écart-type.
  \item On peut visualiser : histogramme, barres, points.
\end{itemize}
\begin{block}{Idée économique}
Deux pays ou deux groupes peuvent avoir le même niveau moyen mais des inégalités ou dispersions très différentes.
\end{block}
\end{frame}

\begin{frame}[fragile]{Écarts à la moyenne}
Pour \code{groupe\_b = [0, 5, 10, 15, 20]}, la moyenne vaut 10.
\begin{lstlisting}
ecart = valeur - moyenne
\end{lstlisting}
\begin{tabular}{@{}rrrrrr@{}}
\toprule
Valeur & 0 & 5 & 10 & 15 & 20 \\
Écart & -10 & -5 & 0 & 5 & 10 \\
\bottomrule
\end{tabular}
\vspace{0.45cm}
\begin{block}{Problème}
Si on additionne les écarts, les valeurs négatives et positives se compensent. Il faut donc mesurer la taille des écarts autrement.
\end{block}
\end{frame}

\begin{frame}{Variance et écart-type : intuition}
\begin{block}{Variance}
La variance est la moyenne des carrés des écarts à la moyenne.
\[
\text{variance} = \frac{1}{n}\sum (x_i - \bar{x})^2
\]
\end{block}
\begin{itemize}
  \item On met les écarts au carré pour éviter les compensations.
  \item Les grands écarts comptent davantage.
  \item L'écart-type est la racine carrée de la variance.
  \item L'écart-type revient dans l'unité de départ, donc il est plus lisible.
\end{itemize}
\end{frame}

\begin{frame}[fragile]{Pseudo-code : variance sans panique}
\begin{lstlisting}
# Objectif : mesurer la dispersion d'une liste
# Entree : liste de valeurs
# 1. Calculer la moyenne
# 2. Pour chaque valeur :
#       calculer l'ecart a la moyenne
#       mettre cet ecart au carre
# 3. Calculer la moyenne des ecarts au carre
# 4. Prendre la racine pour obtenir l'ecart-type
# Sortie : variance et ecart-type
\end{lstlisting}
\begin{block}{Position dans le cours}
À connaître intuitivement. En obligatoire, on insiste surtout sur moyenne, médiane, min, max et graphiques.
\end{block}
\end{frame}

\begin{frame}[fragile]{Livecoding 1 : résumer deux groupes}
Départ en direct :
\begin{lstlisting}
groupe_a = [8, 9, 10, 11, 12]
groupe_b = [0, 5, 10, 15, 20]

def moyenne(valeurs):
    return sum(valeurs) / len(valeurs)

print(moyenne(groupe_a))
print(moyenne(groupe_b))
\end{lstlisting}
À coder ensemble : min, max, étendue, puis une phrase d'interprétation.
\end{frame}

\begin{frame}[fragile]{Livecoding 2 : médiane et dispersion}
Objectif : transformer une idée statistique en fonction Python.
\begin{lstlisting}
def mediane_impair(valeurs):
    valeurs_triees = sorted(valeurs)
    milieu = len(valeurs_triees) // 2
    return valeurs_triees[milieu]

def etendue(valeurs):
    return max(valeurs) - min(valeurs)
\end{lstlisting}
Question en direct : quel groupe a la plus grande étendue ? Est-ce visible dans la moyenne ?
\end{frame}

\begin{frame}[fragile]{Livecoding 3 : variance pas à pas}
Version volontairement explicite, pour comprendre la logique.
\begin{lstlisting}
def variance(valeurs):
    m = moyenne(valeurs)
    carres = []
    for x in valeurs:
        ecart = x - m
        carres.append(ecart ** 2)
    return moyenne(carres)
\end{lstlisting}
Extension en direct si le groupe suit :
\begin{lstlisting}
import math
print(math.sqrt(variance(groupe_b)))
\end{lstlisting}
\end{frame}

\begin{frame}{Graphiques : voir ce que les chiffres cachent}
\begin{tabular}{@{}lll@{}}
\toprule
Graphique & Sert à voir & TP \\
\midrule
Barres & comparer des catégories & TP4, TP5 \\
Histogramme & forme d'une distribution & TP4 \\
Courbe & évolution dans le temps & extension TP5 \\
Nuage de points & relation entre deux variables & plus tard \\
\bottomrule
\end{tabular}
\vspace{0.45cm}
\begin{block}{Règle}
Un graphique doit répondre à une question précise. Il n'est pas une décoration.
\end{block}
\end{frame}

\begin{frame}[fragile]{Livecoding 4 : visualiser une distribution}
\begin{lstlisting}
import matplotlib.pyplot as plt

valeurs = [46, 54, 34, 38, 22, 61]
pays = ["France", "Germany", "Spain", "Italy", "Poland", "Sweden"]

plt.bar(pays, valeurs)
plt.xticks(rotation=45, ha="right")
plt.ylabel("PIB par habitant, milliers de dollars")
plt.tight_layout()
plt.show()
\end{lstlisting}
À demander : quel pays tire la moyenne vers le haut ?
\end{frame}

\begin{frame}{Lire un fichier Our World in Data}
\begin{tabular}{@{}lll@{}}
\toprule
Colonne OWID & Sens & Exemple \\
\midrule
\code{Entity} & pays ou région & France \\
\code{Code} & code pays & FRA \\
\code{Year} & année & 2022 \\
Dernière colonne & indicateur & PIB par habitant \\
\bottomrule
\end{tabular}
\vspace{0.45cm}
\begin{block}{Astuce}
Les fichiers OWID ont souvent une structure régulière : identifier la colonne de valeur est la première étape.
\end{block}
\end{frame}

\begin{frame}[fragile]{Pseudo-code : pipeline pandas OWID}
\begin{lstlisting}
# Objectif : decrire un indicateur OWID
# 1. Lire le CSV depuis une URL
# 2. Afficher les premieres lignes
# 3. Renommer Entity, Code, Year, valeur
# 4. Filtrer une annee recente
# 5. Selectionner quelques pays
# 6. Calculer moyenne, mediane, min, max
# 7. Tracer un graphique
# 8. Rediger une interpretation prudente
\end{lstlisting}
\end{frame}

\begin{frame}[fragile]{Livecoding 5 : lire un CSV OWID}
\begin{lstlisting}
import pandas as pd

url = "https://ourworldindata.org/grapher/gdp-per-capita-worldbank.csv"
df = pd.read_csv(url)

print(df.head())
print(df.shape)
print(df.columns)
\end{lstlisting}
Question : quelle colonne contient l'indicateur économique ?
\end{frame}

\begin{frame}[fragile]{Livecoding 6 : nettoyer et filtrer}
\begin{lstlisting}
value_col = df.columns[-1]
df = df.rename(columns={
    "Entity": "country",
    "Code": "code",
    "Year": "year",
    value_col: "gdp_pc"
})

latest_year = df["year"].max()
latest = df.loc[df["year"] == latest_year].dropna(subset=["gdp_pc"])
\end{lstlisting}
À demander : pourquoi retirer les valeurs manquantes avant de calculer ?
\end{frame}

\begin{frame}[fragile]{Livecoding 7 : calculer et comparer}
\begin{lstlisting}
pays = ["France", "Germany", "Spain", "Poland", "United States", "India"]
selection = latest.loc[latest["country"].isin(pays)].copy()

print(selection[["country", "gdp_pc"]])
print(selection["gdp_pc"].mean())
print(selection["gdp_pc"].median())
print(selection["gdp_pc"].min())
print(selection["gdp_pc"].max())
\end{lstlisting}
Interprétation : la France est-elle au-dessus ou au-dessous de la moyenne de cette sélection ?
\end{frame}

\begin{frame}{Décrire sans surinterpréter}
\begin{block}{Phrase acceptable}
Dans cette sélection de pays et pour cette année, la France est au-dessus ou au-dessous de la moyenne observée du PIB par habitant.
\end{block}
\begin{block}{Phrase trop forte}
La France est plus riche parce que sa politique économique est meilleure.
\end{block}
\begin{itemize}
  \item On décrit des données.
  \item On ne prouve pas une cause.
  \item On précise toujours le périmètre : pays, année, indicateur, source.
\end{itemize}
\end{frame}

\begin{frame}{Ce que les TPs rendent obligatoire}
\begin{tabular}{@{}ll@{}}
\toprule
TP & Obligatoire \\
\midrule
TP3 & fonctions, moyenne, médiane, min, max, comparaison \\
TP4 & graphique, simulation simple, distribution \\
TP5 & CSV OWID, \code{pandas}, graphique, interprétation \\
\bottomrule
\end{tabular}
\vspace{0.45cm}
\begin{block}{Priorité pédagogique}
Écrire du Python clair avant d'empiler les notions statistiques.
\end{block}
\end{frame}

\begin{frame}{Résumé}
\begin{itemize}
  \item Moyenne et médiane résument le centre d'une distribution.
  \item Elles ne captent pas toute la dispersion.
  \item Variance et écart-type mesurent les écarts à la moyenne, mais restent secondaires dans ce cours.
  \item Les graphiques aident à voir ce qu'un chiffre unique cache.
  \item \code{pandas} permet de passer d'une petite liste à un vrai tableau OWID.
  \item Toute interprétation doit rester descriptive et prudente.
\end{itemize}
\end{frame}

\end{document}
